%!TEX root = ../main.tex

\addchap{Abkürzungsverzeichnis}

\begin{acronym} %[ABCDEF]               % längste Abkürzung
	%\setlength{\itemsep}{-\parsep}     % Kein Abstand, kompakte Darstellung
	% Sortiere von Hand oder automatisch mit Kommandozeile (windows): sort file.txt /O file.txt

	\acro{ATLAS}{ATLAS-Laserapplikationsdatenbank}
	\acro{API}{Application Programming Interface (deutsch: Programmierschnittstelle)}
	\acro{ASP}{Active Server Pages}
	\acro{CRUD}{Create, Read, Update, Delete (deutsch: Erstellen, Lesen, Aktualisieren, Löschen)}
	\acro{DHBW}{Dualen Hochschule Baden-Württemberg}
	\acro{HTTP}{Hypertext Transfer Protocol}
	\acro{HTTPS}{Hypertext Transfer Protocol Secure}
	\acro{JSON}{JavaScript Object Notation}
	\acro{JWT}{JSON Web Token}
	\acro{LAC}{Laser Application Center - Abteilung bei TRUMPF}
	\acro{NET}{.NET-Entwicklungsplattform von Microsoft}
	\acro{SPA}{Single-Page-Application (deutsch: Einzelseitenanwendung)}
	\acro{UI}{User Interface (deutsch: Benutzeroberfläche)}
	\acro{VM}{Virtuelle Maschine}

\end{acronym}
